% XPU graph 研究（三卡）：图生文解码、扩散步、TP2 图捕获、去噪步 kernel 分布
\setpart{BAGEL 实测 · XPU graph 与 kernel 分布（三卡）}

\begin{frame}[t]{XPU graph（三卡，卡 0、1、2）：图生文解码，中位数，秒}
\scriptsize\def\codesize{\relax}
\renewcommand{\arraystretch}{1.12}\setlength{\tabcolsep}{3pt}
% generated by scripts/make_tex_tables.py -- do not edit
\scriptsize
\begin{tabularx}{\textwidth}{>{\bfseries}p{4.6cm} c r r r r r}
\toprule
\rowcolor{Mist}
配置 & 缓存 & I2T 1 token eager & graph & I2T 64 token eager & graph & 64 token 变化 \\
\midrule
M* CFG 并行（三卡） & 开 & 0.089 & 0.085 & 2.42 & 1.95 & $-$20\% \\
M* CFG 并行（三卡） & 关 & 0.304 & 0.300 & 2.68 & 2.15 & $-$20\% \\
M* TP2（三卡，LLM 卡 1、2） & 开 & 0.258 & 0.226 & 3.05 & 1.52 & $-$50\% \\
vLLM-Omni two-stage（三卡，Thinker TP2 卡 0、1） & 关 & 0.127 & 0.125 & 1.36 & 1.10 & $-$19\% \\
vLLM-Omni two-stage（三卡，Thinker TP2 卡 0、1） & 开 & 0.054 & 0.051 & 1.25 & 1.02 & $-$19\% \\
\bottomrule
\end{tabularx}

\vspace{4pt}
\begin{itemize}\scriptsize\setlength\itemsep{1pt}
  \item 配置：M* CFG 并行，三个 rank 分别为卡 0 的主分支、卡 1 与卡 2 的两路引导分支；M* TP2，卡 0 为 ViT 与 VAE，卡 1、2 切分 LLM；Omni two-stage，Thinker TP2 在卡 0、1，DiT 在卡 2（图生文不使用 DiT）。
  \item graph 只作用于 LLM 解码步。M* CFG 并行的 graph 无需补丁；M* TP2 的 graph 需要补丁（见下页）。
  \item 两边图像输入规模一致：M* 为默认预处理的 1369 个图像 patch，Omni 提示为 1386 token（含 17 个文本 token）。
\end{itemize}
\takeaway{三卡下 graph 使 M* 的 64 token 解码在 CFG 并行中快 \pct{gCfgOn64}，在 TP2 中快 \pct{gTp2G64}。M* TP2 graph 后为 \val{tp2g/mstar-tp2-graph-fixed/i2t64} s，Omni two-stage graph 为 \val{g3/omni-tp2-graph-chat/i2t64} s，M* 仍慢 \pct{gTp2Gap}。}
\end{frame}

\begin{frame}[t]{XPU graph（三卡）：扩散步不变，TP2 图捕获的限制}
\scriptsize\def\codesize{\relax}
\renewcommand{\arraystretch}{1.12}\setlength{\tabcolsep}{3pt}
% generated by scripts/make_tex_tables.py -- do not edit
\scriptsize
\begin{tabularx}{\textwidth}{>{\bfseries}p{5.6cm} r r r r r}
\toprule
\rowcolor{Mist}
配置 & T2I eager & T2I graph & I2I eager & I2I graph & I2I 变化 \\
\midrule
M* CFG 并行（三卡） & 7.73 & 7.75 & 9.28 & 9.26 & $-$0\% \\
M* TP2（三卡，LLM 卡 1、2） & 7.68 & 7.64 & -- & -- & n/a \\
vLLM-Omni 单 stage CFG 并行（三卡） & -- & -- & 11.19 & -- & n/a \\
vLLM-Omni two-stage（三卡，Thinker TP2 卡 0、1，DiT 卡 2） & 17.91 & 18.01 & -- & -- & n/a \\
\bottomrule
\end{tabularx}

\vspace{4pt}
\begin{itemize}\scriptsize\setlength\itemsep{1pt}
  \item 扩散步两边都是 eager，graph 不作用于扩散。M* CFG 并行的 T2I 与 I2I 在 graph 前后基本持平（T2I \pct{gDiffT}，I2I \pct{gDiffI}）。vLLM-Omni 的 DiT 不支持 denoise graph。
  \item M* TP2 图捕获失败的原因有三：(1) XCCL 的 \code{all\_reduce} 消息超过约 4 MiB 时无法进入 graph；(2) inductor 生成的 functional \code{all\_gather}（lm\_head 汇聚 logits）在 XPU graph 中不可用，普通 \code{all\_gather} 在约 4 行 logits 以内可用；(3) 最大的 prefill 捕获失败后，捕获状态残留，decode 随之失败。
  \item 补丁（\code{scripts/patches/mstar\_xpu\_tp2\_graph.patch}）作用于已安装的 M* 包，未合入上游：\code{all\_gather} 不参与 compile；prefill 捕获 token 上限 512；decode 与 prefill 的捕获 batch size 限为 [1, 2, 4]。
  \item 两卡 CFG 并行的 graph 因显存不足无法启动（每卡约 27 GiB 的完整模型）。
\end{itemize}
\takeaway{graph 只加速图生文的解码：64 token 快 19--50\%；1 token（预填充为主）只快 1--12\%。文生图与图像编辑的扩散步不经过 graph，前后差异在 1\% 以内。M* TP2 的 graph 需要补丁，CFG 并行则不需要。}
\end{frame}

\begin{frame}[t]{去噪步的 kernel 分布（单步，rank 0）}
\scriptsize\def\codesize{\relax}
\renewcommand{\arraystretch}{1.08}\setlength{\tabcolsep}{3pt}
% generated by scripts/make_kernel_table.py -- do not edit
\scriptsize\setlength{\tabcolsep}{2.5pt}
\begin{tabularx}{\textwidth}{>{\bfseries}p{5.0cm} r r r r r}
\toprule
\rowcolor{Mist}
kernel 类别 & Omni ms & Omni 占比 & Omni 次数 & M* 占比 & M* 次数 \\
\midrule
gemm（矩阵乘） & 119.7 & 61.8\% & 229 & 84.1\% & 78 \\
attention（FMHA） & 5.9 & 3.1\% & 28 & 4.5\% & 18 \\
融合 kernel（Triton / vLLM） & 8.8 & 4.5\% & 56 & 10.5\% & 351 \\
index、gather、scatter & 27.1 & 14.0\% & 1017 & 0.8\% & 18 \\
copy、cat & 15.4 & 8.0\% & 1021 & 0.0\% & 0 \\
RMSNorm 拆散部分（pow、mean、reduce） & 5.7 & 2.9\% & 452 & 0.0\% & 0 \\
其他 elementwise & 11.2 & 5.8\% & 1396 & 0.1\% & 18 \\
\midrule
合计（kernel 忙碌时间） & 193.7 & 100\% & 4199 & 100\% & 483 \\
窗口时长 & 201.7 & & & 57.6 & \\
空闲（窗口 − 忙碌） & 8.0 & 4.0\% & & 11.2\% & \\
gemm 平均耗时（ms） & 0.52 & & & 0.55 & \\
\bottomrule
\end{tabularx}

\vspace{4pt}
\begin{itemize}\scriptsize\setlength\itemsep{1pt}
  \item Omni：单 stage CFG 并行，卡 0、1（完整一步）。M*：CFG 并行，卡 0、1、2（58 ms 片段，只比较占比）。
\end{itemize}
\takeaway{Omni 的非 gemm、非 attention 小 kernel 约 68 ms（占 35\%）；单个 gemm 两边平均耗时相同（0.52 与 0.55 ms）。M* 整步未取得，整步差距未完全分配。}
\end{frame}
